\section{Problem Definition}
\label{sec:problem}

We model co-evolutionary agent security as a partially observable,
repeated interaction between finite attacker and defender populations
$\mathcal{A}_t$ and $\mathcal{D}_t$ over episodes $t=1,\ldots,T$. Each
episode consists of up to $H$ turns in a cyber environment $\mathcal{E}$.
Agents do not observe the full state, the internal policies of others, or
future environment draws. Partial observability is essential: with full
information the interaction collapses into a complete-information game
ill-suited to modelling operational uncertainty.

An agent is defined as
\begin{equation}
  A_i \;=\; \bigl(M_i,\, G_i,\, \mathcal{T}_i,\, \mathcal{C}_i,\,
  K_i,\, \pi_i,\, \mu_i\bigr),
  \label{eq:agent}
\end{equation}
where $M_i$ is the underlying model (heuristic or LLM), $G_i$ a
high-level objective, $\mathcal{T}_i$ the tool set exposed by the
environment, $\mathcal{C}_i$ implied capabilities, $K_i$ revealed
knowledge, $\pi_i$ the policy, and $\mu_i$ memory. Table~\ref{tab:notation}
summarises the main symbols used in later sections.

\begin{table}[!t]
\caption{Principal notation.}
\label{tab:notation}
\centering
\setlength{\tabcolsep}{3pt}
\small
\begin{tabular}{@{}ll@{}}
\toprule
\textbf{Symbol} & \textbf{Meaning} \\
\midrule
$\mathcal{A}_t$, $\mathcal{D}_t$ & Attacker / defender populations \\
$\mathcal{E}$, $H$, $T$ & Environment, turn budget, episodes \\
L0, L1, L2 & Static, episodic, persistent adaptation \\
$s_i$, CEP & Strategy signature, co-evolutionary pressure \\
ASR, DR, ESR & Attack success, detection, emergent risk \\
$B_i^{\mathrm{enc}}$ & Encoded behaviour set for role $i$ \\
\bottomrule
\end{tabular}
\end{table}

The experimenter
sets $G_i$ and the adaptation regime; the environment sets
$\mathcal{T}_i$; $\pi_i$ is \emph{not} a hard-coded attack or defence
script~\cite{yao2023react}. Separating capability, objective and
strategy ensures that a surprising tool sequence can be analysed as
strategy discovery rather than as a newly granted power.

Adaptation is classified into three levels used throughout the paper:
\textbf{L0 (static)} leaves $\pi_i$ and $\mu_i$ fixed across episodes;
\textbf{L1 (episodic)} allows within-episode learning but wipes memory at
the episode boundary; \textbf{L2 (persistent)} retains a compact
cross-episode reflection in $\mu_i$. Populations may communicate through a
directed graph $G_t=(V_t,E_t)$~\cite{wu2023autogen,li2023camel}; coalitions are
detected \emph{post hoc} from reciprocal same-role traffic and are never
objectives in themselves.

Table~\ref{tab:rq} lists the research questions. RQ0--RQ1, RQ3 and RQ5 are
evaluated in Section~\ref{sec:results} on one-versus-one cells; RQ2 (emergence
vs.\ $B_i^{\mathrm{enc}}$) is deferred because ladder pacing behaviours are
declared in $B_A^{\mathrm{enc}}$; RQ4 (population scaling) is left for future
work.

\begin{table}[!t]
\caption{Research questions and their operationalisation.}
\label{tab:rq}
\centering
\setlength{\tabcolsep}{3pt}
\small
\begin{tabularx}{\columnwidth}{@{}lX@{}}
\toprule
\textbf{RQ} & \textbf{Question / metric} \\
\midrule
RQ0 & How does security behaviour evolve under repeated interaction? \\
RQ1 & Effect of adaptation on ASR and DR (ladder factorial) \\
RQ2 & Emergent behaviours vs.\ $B_i^{\mathrm{enc}}$ (deferred) \\
RQ3 & Co-evolution measured by CEP and strategy signatures \\
RQ4 & Population size, topology, communication (E4--E8; future) \\
RQ5 & Dynamics class: equilibrium, convergence, non-convergent change \\
\bottomrule
\end{tabularx}
\end{table}

Each turn proceeds in fixed order (attacker then defender) with environment-mediated
feedback after each action; neither side observes the other's internal state. Episode termination
occurs when the attacker exfiltrates, the defender contains all footholds,
or the turn budget expires. Between episodes, only persistent memory and
adaptive parameters carry forward; the environment is re-sampled from the
same generative process with an incremented seed so that cross-episode
learning cannot memorise a single topology instance.
